\documentclass[12pt]{article}

\usepackage{amsmath}
\usepackage{amssymb}
\usepackage[margin=0.6in]{geometry}

\setlength{\abovedisplayskip}{4pt}
\setlength{\belowdisplayskip}{4pt}
\setlength{\abovedisplayshortskip}{2pt}
\setlength{\belowdisplayshortskip}{2pt}
\setlength{\parskip}{2pt}
\setlength{\parsep}{0pt}

\begin{document}
	
	\begin{center}
		{\large\textbf{APM1220}} \quad
		\textbf{FA 1 -- Foundations: Matrix Algebra} \quad 
		\\ Joseph Raphael Espiritu
	\end{center}
	
	\vspace{-0.3em}
	
	\section*{Problem 2.2}
	
	Given
	
	\[
	A =
	\begin{bmatrix}
		-1 & 3\\
		4 & 2
	\end{bmatrix},
	\qquad
	B =
	\begin{bmatrix}
		4 & -3\\
		1 & -2\\
		-2 & 0
	\end{bmatrix},
	\qquad
	C =
	\begin{bmatrix}
		5\\
		-4\\
		2
	\end{bmatrix}.
	\]
	
	\subsection*{(a) $5A$}
	
	Scalar multiplication is performed by multiplying every element of $A$
	by $5$:
	
	\[
	5A
	=
	5
	\begin{bmatrix}
		-1 & 3\\
		4 & 2
	\end{bmatrix}
	\]
	
	\[
	=
	\begin{bmatrix}
		5(-1) & 5(3)\\
		5(4) & 5(2)
	\end{bmatrix}
	\]
	
	\[
	=
	\begin{bmatrix}
		-5 & 15\\
		20 & 10
	\end{bmatrix}.
	\]
	
	Therefore,
	
	\[
	\boxed{
		5A =
		\begin{bmatrix}
			-5 & 15\\
			20 & 10
	\end{bmatrix}}
	\]
	
	\subsection*{(b) $BA$}
	
	The dimensions are
	
	\[
	B_{3\times2}A_{2\times2}.
	\]
	
	Since the inner dimensions are equal,
	
	\[
	(3\times2)(2\times2)=(3\times2),
	\]
	
	so $BA$ is defined.
	
	\[
	BA
	=
	\begin{bmatrix}
		4 & -3\\
		1 & -2\\
		-2 & 0
	\end{bmatrix}
	\begin{bmatrix}
		-1 & 3\\
		4 & 2
	\end{bmatrix}
	\]
	
	\[
	=
	\begin{bmatrix}
		4(-1)+(-3)(4) &
		4(3)+(-3)(2)
		\\
		1(-1)+(-2)(4) &
		1(3)+(-2)(2)
		\\
		(-2)(-1)+0(4) &
		(-2)(3)+0(2)
	\end{bmatrix}
	\]
	
	\[
	=
	\begin{bmatrix}
		-4-12 & 12-6\\
		-1-8 & 3-4\\
		2+0 & -6+0
	\end{bmatrix}
	\]
	
	\[
	=
	\begin{bmatrix}
		-16 & 6\\
		-9 & -1\\
		2 & -6
	\end{bmatrix}.
	\]
	
	Therefore,
	
	\[
	\boxed{
		BA =
		\begin{bmatrix}
			-16 & 6\\
			-9 & -1\\
			2 & -6
	\end{bmatrix}}
	\]
	
	\subsection*{(c) $A'B'$}
	
	First, transpose $A$ and $B$:
	
	\[
	A'
	=
	\begin{bmatrix}
		-1 & 4\\
		3 & 2
	\end{bmatrix}
	\]
	
	and
	
	\[
	B'
	=
	\begin{bmatrix}
		4 & 1 & -2\\
		-3 & -2 & 0
	\end{bmatrix}.
	\]
	
	The dimensions are
	
	\[
	A'_{2\times2}B'_{2\times3}.
	\]
	
	Therefore,
	
	\[
	(2\times2)(2\times3)=(2\times3),
	\]
	
	so $A'B'$ is defined.
	
	\[
	A'B'
	=
	\begin{bmatrix}
		-1 & 4\\
		3 & 2
	\end{bmatrix}
	\begin{bmatrix}
		4 & 1 & -2\\
		-3 & -2 & 0
	\end{bmatrix}
	\]
	
	\[
	=
	\begin{bmatrix}
		(-1)(4)+4(-3) &
		(-1)(1)+4(-2) &
		(-1)(-2)+4(0)
		\\
		3(4)+2(-3) &
		3(1)+2(-2) &
		3(-2)+2(0)
	\end{bmatrix}
	\]
	
	\[
	=
	\begin{bmatrix}
		-4-12 & -1-8 & 2+0\\
		12-6 & 3-4 & -6+0
	\end{bmatrix}
	\]
	
	\[
	=
	\begin{bmatrix}
		-16 & -9 & 2\\
		6 & -1 & -6
	\end{bmatrix}.
	\]
	
	Therefore,
	
	\[
	\boxed{
		A'B' =
		\begin{bmatrix}
			-16 & -9 & 2\\
			6 & -1 & -6
	\end{bmatrix}}
	\]
	
	\subsection*{(d) $C'B$}
	
	First, transpose $C$:
	
	\[
	C'
	=
	\begin{bmatrix}
		5 & -4 & 2
	\end{bmatrix}.
	\]
	
	The dimensions are
	
	\[
	C'_{1\times3}B_{3\times2}.
	\]
	
	Therefore,
	
	\[
	(1\times3)(3\times2)=(1\times2),
	\]
	
	so $C'B$ is defined.
	
	\[
	C'B
	=
	\begin{bmatrix}
		5 & -4 & 2
	\end{bmatrix}
	\begin{bmatrix}
		4 & -3\\
		1 & -2\\
		-2 & 0
	\end{bmatrix}
	\]
	
	\[
	=
	\begin{bmatrix}
		5(4)+(-4)(1)+2(-2)
		&
		5(-3)+(-4)(-2)+2(0)
	\end{bmatrix}
	\]
	
	\[
	=
	\begin{bmatrix}
		20-4-4 &
		-15+8+0
	\end{bmatrix}
	\]
	
	\[
	=
	\begin{bmatrix}
		12 & -7
	\end{bmatrix}.
	\]
	
	Therefore,
	
	\[
	\boxed{
		C'B =
		\begin{bmatrix}
			12 & -7
	\end{bmatrix}}
	\]
	
	\subsection*{(e) Is $AB$ defined?}
	
	The dimensions are
	
	\[
	A_{2\times2}B_{3\times2}.
	\]
	
	For matrix multiplication to be defined, the number of columns
	of the first matrix must equal the number of rows of the second matrix.
	
	Here,
	
	\[
	2 \neq 3.
	\]
	
	Therefore,
	
	\[
	\boxed{AB\text{ is not defined.}}
	\]
	
	\section*{Problem 2.4}
	
	When $A^{-1}$ and $B^{-1}$ exist, prove each of the following:
	
	\subsection*{(a) $(A')^{-1}=(A^{-1})'$}
	
	By the definition of a matrix inverse,
	
	\[
	AA^{-1}=A^{-1}A=I.
	\]
	
	We first transpose the equation
	
	\[
	AA^{-1}=I.
	\]
	
	Taking the transpose of both sides gives
	
	\[
	(AA^{-1})'=I'.
	\]
	
	The identity matrix is unchanged by transposition because its
	nonzero entries lie only on the main diagonal. Therefore,
	
	\[
	I'=I.
	\]
	
	Also, the transpose of a matrix product reverses the order:
	
	\[
	(AB)'=B'A'.
	\]
	
	Thus,
	
	\[
	(AA^{-1})'
	=
	(A^{-1})'A'.
	\]
	
	Therefore,
	
	\[
	(A^{-1})'A'=I.
	\]
	
	We also start with the other inverse relationship,
	
	\[
	A^{-1}A=I,
	\]
	
	and transpose both sides:
	
	\[
	(A^{-1}A)'=I'.
	\]
	
	Using the transpose product rule,
	
	\[
	A'(A^{-1})'=I.
	\]
	
	Hence, we have shown both
	
	\[
	(A^{-1})'A'=I
	\]
	
	and
	
	\[
	A'(A^{-1})'=I.
	\]
	
	Therefore, $(A^{-1})'$ satisfies the definition of the inverse of $A'$.
	Thus,
	
	\[
	\boxed{(A')^{-1}=(A^{-1})'}.
	\]
	
	\subsection*{(b) $(AB)^{-1}=B^{-1}A^{-1}$}
	
	We want to prove that
	
	\[
	\boxed{(AB)^{-1}=B^{-1}A^{-1}}.
	\]
	
	Rather than assuming that the equality is already true, we verify
	whether $B^{-1}A^{-1}$ satisfies the definition of the inverse of $AB$.
	
	For a matrix $C$ to be the inverse of $AB$, it must satisfy
	
	\[
	C(AB)=(AB)C=I.
	\]
	
	We propose
	
	\[
	C=B^{-1}A^{-1}.
	\]
	
	We first verify the product in one direction:
	
	\[
	(B^{-1}A^{-1})(AB).
	\]
	
	Using associativity of matrix multiplication, we regroup the factors:
	
	\[
	B^{-1}(A^{-1}A)B.
	\]
	
	Since $A^{-1}$ is the inverse of $A$,
	
	\[
	A^{-1}A=I.
	\]
	
	Therefore,
	
	\[
	B^{-1}IB.
	\]
	
	The identity matrix does not change a matrix when multiplied by it:
	
	\[
	B^{-1}B.
	\]
	
	Since $B^{-1}$ is the inverse of $B$,
	
	\[
	B^{-1}B=I.
	\]
	
	Hence,
	
	\[
	\boxed{(B^{-1}A^{-1})(AB)=I}.
	\]
	
	We must also verify the product in the other direction:
	
	\[
	(AB)(B^{-1}A^{-1}).
	\]
	
	Again, regrouping the factors gives
	
	\[
	A(BB^{-1})A^{-1}.
	\]
	
	Since
	
	\[
	BB^{-1}=I,
	\]
	
	we obtain
	
	\[
	AIA^{-1}.
	\]
	
	Using the identity property,
	
	\[
	AA^{-1}.
	\]
	
	Since
	
	\[
	AA^{-1}=I,
	\]
	
	we obtain
	
	\[
	\boxed{(AB)(B^{-1}A^{-1})=I}.
	\]
	
	Therefore, both inverse conditions have been satisfied:
	
	\[
	(B^{-1}A^{-1})(AB)=I
	\]
	
	and
	
	\[
	(AB)(B^{-1}A^{-1})=I.
	\]
	
	Thus, $B^{-1}A^{-1}$ is the inverse of $AB$. Therefore,
	
	\[
	\boxed{(AB)^{-1}=B^{-1}A^{-1}}.
	\]
	
	\section*{Problem 2.8}
	
	Given
	
	\[
	A=
	\begin{bmatrix}
		1&2\\
		2&-2
	\end{bmatrix}.
	\]
	
	We find the eigenvalues, corresponding normalized eigenvectors,
	and the spectral decomposition of $A$.
	
	\subsection*{1. Eigenvalues}
	
	An eigenvector $x\neq0$ satisfies
	
	\[
	Ax=\lambda x.
	\]
	
	Since $Ix=x$, we can write
	
	\[
	Ax=\lambda Ix.
	\]
	
	Therefore,
	
	\[
	(A-\lambda I)x=0.
	\]
	
	For a nonzero solution $x$ to exist, $A-\lambda I$ must be singular.
	Thus, the characteristic equation is
	
	\[
	\det(A-\lambda I)=0.
	\]
	
	The identity matrix is
	
	\[
	I=
	\begin{bmatrix}
		1&0\\
		0&1
	\end{bmatrix},
	\]
	
	so
	
	\[
	A-\lambda I
	=
	\begin{bmatrix}
		1-\lambda&2\\
		2&-2-\lambda
	\end{bmatrix}.
	\]
	
	Therefore,
	
	\[
	\begin{aligned}
		\det(A-\lambda I)
		&=
		\begin{vmatrix}
			1-\lambda&2\\
			2&-2-\lambda
		\end{vmatrix}\\
		&=(1-\lambda)(-2-\lambda)-(2)(2)\\
		&=\lambda^2+\lambda-6.
	\end{aligned}
	\]
	
	Set the determinant equal to zero:
	
	\[
	\lambda^2+\lambda-6=0.
	\]
	
	Factoring gives
	
	\[
	(\lambda-2)(\lambda+3)=0.
	\]
	
	Hence, the eigenvalues are
	
	\[
	\boxed{\lambda_1=2,\qquad\lambda_2=-3}.
	\]
	
	\subsection*{2. Eigenvector for $\lambda_1=2$}
	
	For $\lambda_1=2$,
	
	\[
	(A-2I)x_1=0.
	\]
	
	We have
	
	\[
	\begin{aligned}
		A-2I
		&=
		\begin{bmatrix}
			1&2\\
			2&-2
		\end{bmatrix}
		-
		\begin{bmatrix}
			2&0\\
			0&2
		\end{bmatrix}\\
		&=
		\begin{bmatrix}
			-1&2\\
			2&-4
		\end{bmatrix}.
	\end{aligned}
	\]
	
	Therefore,
	
	\[
	\begin{bmatrix}
		-1&2\\
		2&-4
	\end{bmatrix}
	\begin{bmatrix}
		x_1\\
		x_2
	\end{bmatrix}
	=
	\begin{bmatrix}
		0\\
		0
	\end{bmatrix}.
	\]
	
	This gives
	
	\[
	-x_1+2x_2=0.
	\]
	
	Hence,
	
	\[
	x_1=2x_2.
	\]
	
	Choosing $x_2=1$ gives
	
	\[
	\boxed{
		x_1=
		\begin{bmatrix}
			2\\
			1
		\end{bmatrix}
	}.
	\]
	
	\subsection*{3. Normalize the first eigenvector}
	
	The length of $x_1$ is
	
	\[
	\begin{aligned}
		\|x_1\|
		&=\sqrt{x_1'x_1}\\
		&=
		\sqrt{
			\begin{bmatrix}
				2&1
			\end{bmatrix}
			\begin{bmatrix}
				2\\
				1
			\end{bmatrix}
		}\\
		&=\sqrt{4+1}\\
		&=\sqrt{5}.
	\end{aligned}
	\]
	
	The normalized eigenvector is
	
	\[
	e_1=\frac{x_1}{\|x_1\|}.
	\]
	
	Thus,
	
	\[
	\boxed{
		e_1=
		\frac{1}{\sqrt{5}}
		\begin{bmatrix}
			2\\
			1
		\end{bmatrix}
	}.
	\]
	
	It satisfies
	
	\[
	e_1'e_1=1.
	\]
	
	\subsection*{4. Eigenvector for $\lambda_2=-3$}
	
	For $\lambda_2=-3$,
	
	\[
	(A+3I)x_2=0.
	\]
	
	We have
	
	\[
	\begin{aligned}
		A+3I
		&=
		\begin{bmatrix}
			1&2\\
			2&-2
		\end{bmatrix}
		+
		\begin{bmatrix}
			3&0\\
			0&3
		\end{bmatrix}\\
		&=
		\begin{bmatrix}
			4&2\\
			2&1
		\end{bmatrix}.
	\end{aligned}
	\]
	
	Therefore,
	
	\[
	\begin{bmatrix}
		4&2\\
		2&1
	\end{bmatrix}
	\begin{bmatrix}
		x_1\\
		x_2
	\end{bmatrix}
	=
	\begin{bmatrix}
		0\\
		0
	\end{bmatrix}.
	\]
	
	From the first equation,
	
	\[
	4x_1+2x_2=0.
	\]
	
	Dividing by $2$ gives
	
	\[
	2x_1+x_2=0.
	\]
	
	Therefore,
	
	\[
	x_2=-2x_1.
	\]
	
	Choosing $x_1=1$ gives
	
	\[
	\boxed{
		x_2=
		\begin{bmatrix}
			1\\
			-2
		\end{bmatrix}
	}.
	\]
	
	\subsection*{5. Normalize the second eigenvector}
	
	The length of $x_2$ is
	
	\[
	\begin{aligned}
		\|x_2\|
		&=\sqrt{x_2'x_2}\\
		&=
		\sqrt{
			\begin{bmatrix}
				1&-2
			\end{bmatrix}
			\begin{bmatrix}
				1\\
				-2
			\end{bmatrix}
		}\\
		&=\sqrt{1+4}\\
		&=\sqrt{5}.
	\end{aligned}
	\]
	
	Therefore,
	
	\[
	e_2=\frac{x_2}{\|x_2\|},
	\]
	
	giving
	
	\[
	\boxed{
		e_2=
		\frac{1}{\sqrt{5}}
		\begin{bmatrix}
			1\\
			-2
		\end{bmatrix}
	}.
	\]
	
	It satisfies
	
	\[
	e_2'e_2=1.
	\]
	
	The two normalized eigenvectors are perpendicular:
	
	\[
	\begin{aligned}
		e_1'e_2
		&=
		\frac{1}{5}
		\begin{bmatrix}
			2&1
		\end{bmatrix}
		\begin{bmatrix}
			1\\
			-2
		\end{bmatrix}\\
		&=\frac{1}{5}(2-2)\\
		&=0.
	\end{aligned}
	\]
	
	Thus,
	
	\[
	\boxed{e_1'e_2=0}.
	\]
	
	Hence, $e_1$ and $e_2$ are orthonormal eigenvectors.
	
	\subsection*{6. Spectral Decomposition}
	
	For a symmetric matrix, the spectral decomposition is
	
	\[
	A=\sum_i\lambda_i e_i e_i'.
	\]
	
	For the two eigenvalues,
	
	\[
	A=\lambda_1e_1e_1'
	+\lambda_2e_2e_2'.
	\]
	
	Substituting the eigenvalues gives
	
	\[
	A=2e_1e_1'-3e_2e_2'.
	\]
	
	First,
	
	\[
	\begin{aligned}
		e_1e_1'
		&=
		\frac{1}{\sqrt{5}}
		\begin{bmatrix}
			2\\
			1
		\end{bmatrix}
		\frac{1}{\sqrt{5}}
		\begin{bmatrix}
			2&1
		\end{bmatrix}\\
		&=
		\frac{1}{5}
		\begin{bmatrix}
			4&2\\
			2&1
		\end{bmatrix}.
	\end{aligned}
	\]
	
	Therefore,
	
	\[
	2e_1e_1'
	=
	\frac{1}{5}
	\begin{bmatrix}
		8&4\\
		4&2
	\end{bmatrix}.
	\]
	
	Next,
	
	\[
	\begin{aligned}
		e_2e_2'
		&=
		\frac{1}{\sqrt{5}}
		\begin{bmatrix}
			1\\
			-2
		\end{bmatrix}
		\frac{1}{\sqrt{5}}
		\begin{bmatrix}
			1&-2
		\end{bmatrix}\\
		&=
		\frac{1}{5}
		\begin{bmatrix}
			1&-2\\
			-2&4
		\end{bmatrix}.
	\end{aligned}
	\]
	
	Therefore,
	
	\[
	-3e_2e_2'
	=
	\frac{1}{5}
	\begin{bmatrix}
		-3&6\\
		6&-12
	\end{bmatrix}.
	\]
	
	Adding the two terms,
	
	\[
	\begin{aligned}
		A
		&=
		\frac{1}{5}
		\begin{bmatrix}
			8&4\\
			4&2
		\end{bmatrix}
		+
		\frac{1}{5}
		\begin{bmatrix}
			-3&6\\
			6&-12
		\end{bmatrix}\\
		&=
		\frac{1}{5}
		\begin{bmatrix}
			5&10\\
			10&-10
		\end{bmatrix}\\
		&=
		\begin{bmatrix}
			1&2\\
			2&-2
		\end{bmatrix}.
	\end{aligned}
	\]
	
	Therefore, the spectral decomposition is
	
	\[
	\boxed{
		A=2e_1e_1'-3e_2e_2'
	}
	\]
	
	where
	
	\[
	\boxed{
		e_1=
		\frac{1}{\sqrt{5}}
		\begin{bmatrix}
			2\\
			1
		\end{bmatrix},
		\qquad
		e_2=
		\frac{1}{\sqrt{5}}
		\begin{bmatrix}
			1\\
			-2
		\end{bmatrix}
	}.
	\]
	
	\section*{Problem 2.11}
	
	Let $A$ be a $p\times p$ diagonal matrix, where
	
	\[
	A=\{a_{ij}\},
	\qquad
	a_{ij}=0\quad\text{for }i\neq j.
	\]
	
	Thus, $A$ has the form
	
	\[
	A=
	\begin{bmatrix}
		a_{11}&0&0&\cdots&0\\
		0&a_{22}&0&\cdots&0\\
		0&0&a_{33}&\cdots&0\\
		\vdots&\vdots&\vdots&\ddots&\vdots\\
		0&0&0&\cdots&a_{pp}
	\end{bmatrix}.
	\]
	
	We want to show that
	
	\[
	\boxed{
		|A|=a_{11}a_{22}\cdots a_{pp}
	}.
	\]
	
	Using cofactor expansion along the first row,
	
	\[
	|A|
	=
	a_{11}A_{11}
	+a_{12}A_{12}
	+\cdots
	+a_{1p}A_{1p}.
	\]
	
	Since $A$ is diagonal,
	
	\[
	a_{12}=a_{13}=\cdots=a_{1p}=0.
	\]
	
	Therefore,
	
	\[
	|A|
	=
	a_{11}A_{11}.
	\]
	
	The cofactor $A_{11}$ is the determinant of the submatrix
	obtained by deleting the first row and first column:
	
	\[
	A_{11}
	=
	\begin{bmatrix}
		a_{22}&0&\cdots&0\\
		0&a_{33}&\cdots&0\\
		\vdots&\vdots&\ddots&\vdots\\
		0&0&\cdots&a_{pp}
	\end{bmatrix}.
	\]
	
	This submatrix is also diagonal. Expanding its determinant along
	its first row gives
	
	\[
	A_{11}=a_{22}A_{22}.
	\]
	
	Therefore,
	
	\[
	|A|
	=
	a_{11}a_{22}A_{22}.
	\]
	
	Repeating the same argument for each successive diagonal submatrix,
	
	\[
	\begin{aligned}
		|A|
		&=a_{11}a_{22}A_{22}\\
		&=a_{11}a_{22}a_{33}A_{33}\\
		&\;\vdots\\
		&=a_{11}a_{22}\cdots a_{pp}.
	\end{aligned}
	\]
	
	Therefore, the determinant of a diagonal matrix is the product
	of its diagonal elements:
	
	\[
	\boxed{
		|A|=a_{11}a_{22}\cdots a_{pp}
	}.
	\]
	
	\section*{Problem 2.24}
	
	Given
	
	\[
	\Sigma=
	\begin{bmatrix}
		4&0&0\\
		0&9&0\\
		0&0&1
	\end{bmatrix}.
	\]
	
	We use the properties of diagonal and triangular matrices,
	matrix inverses, and eigenvalues.
	
	\subsection*{(a) Find $\Sigma^{-1}$}
	
	Since $\Sigma$ is diagonal, it is also triangular.
	
	For a diagonal matrix with nonzero diagonal elements, its inverse
	is obtained by taking the reciprocal of each diagonal element.
	This follows from
	
	\[
	\Sigma\Sigma^{-1}=I.
	\]
	
	Therefore,
	
	\[
	\boxed{
		\Sigma^{-1}
		=
		\begin{bmatrix}
			1/4&0&0\\
			0&1/9&0\\
			0&0&1
		\end{bmatrix}
	}.
	\]
	
	To verify,
	
	\[
	\begin{aligned}
		\Sigma\Sigma^{-1}
		&=
		\begin{bmatrix}
			4&0&0\\
			0&9&0\\
			0&0&1
		\end{bmatrix}
		\begin{bmatrix}
			1/4&0&0\\
			0&1/9&0\\
			0&0&1
		\end{bmatrix}\\
		&=
		\begin{bmatrix}
			4(1/4)&0&0\\
			0&9(1/9)&0\\
			0&0&1(1)
		\end{bmatrix}\\
		&=
		\begin{bmatrix}
			1&0&0\\
			0&1&0\\
			0&0&1
		\end{bmatrix}\\
		&=I.
	\end{aligned}
	\]
	
	Thus, the matrix above is $\Sigma^{-1}$.
	
	\subsection*{(b) Eigenvalues and Eigenvectors of $\Sigma$}
	
	A useful property of triangular matrices is that their eigenvalues
	are the diagonal elements. Since $\Sigma$ is diagonal,
	
	\[
	\boxed{
		\lambda_1=4,\qquad
		\lambda_2=9,\qquad
		\lambda_3=1
	}.
	\]
	
	To find the corresponding eigenvectors, we use
	
	\[
	\Sigma e_i=\lambda_i e_i.
	\]
	
	Because $\Sigma$ is diagonal, multiplication by $\Sigma$ scales
	each standard basis vector by the corresponding diagonal element.
	Therefore, the standard basis vectors naturally give the
	corresponding eigenvectors.
	
	For $\lambda_1=4$, choose
	
	\[
	e_1=
	\begin{bmatrix}
		1\\0\\0
	\end{bmatrix}.
	\]
	
	This is not an arbitrary assumption. It is the first standard basis
	vector, which isolates the first diagonal element of $\Sigma$:
	
	\[
	\begin{aligned}
		\Sigma e_1
		&=
		\begin{bmatrix}
			4&0&0\\
			0&9&0\\
			0&0&1
		\end{bmatrix}
		\begin{bmatrix}
			1\\0\\0
		\end{bmatrix}\\
		&=
		\begin{bmatrix}
			4\\0\\0
		\end{bmatrix}\\
		&=
		4
		\begin{bmatrix}
			1\\0\\0
		\end{bmatrix}\\
		&=4e_1.
	\end{aligned}
	\]
	
	Therefore,
	
	\[
	\boxed{
		\lambda_1=4,\qquad
		e_1=
		\begin{bmatrix}
			1\\0\\0
		\end{bmatrix}
	}.
	\]
	
	For $\lambda_2=9$, use the second standard basis vector:
	
	\[
	e_2=
	\begin{bmatrix}
		0\\1\\0
	\end{bmatrix}.
	\]
	
	Then,
	
	\[
	\begin{aligned}
		\Sigma e_2
		&=
		\begin{bmatrix}
			4&0&0\\
			0&9&0\\
			0&0&1
		\end{bmatrix}
		\begin{bmatrix}
			0\\1\\0
		\end{bmatrix}\\
		&=
		\begin{bmatrix}
			0\\9\\0
		\end{bmatrix}\\
		&=
		9
		\begin{bmatrix}
			0\\1\\0
		\end{bmatrix}\\
		&=9e_2.
	\end{aligned}
	\]
	
	Therefore,
	
	\[
	\boxed{
		\lambda_2=9,\qquad
		e_2=
		\begin{bmatrix}
			0\\1\\0
		\end{bmatrix}
	}.
	\]
	
	For $\lambda_3=1$, use the third standard basis vector:
	
	\[
	e_3=
	\begin{bmatrix}
		0\\0\\1
	\end{bmatrix}.
	\]
	
	Then,
	
	\[
	\begin{aligned}
		\Sigma e_3
		&=
		\begin{bmatrix}
			4&0&0\\
			0&9&0\\
			0&0&1
		\end{bmatrix}
		\begin{bmatrix}
			0\\0\\1
		\end{bmatrix}\\
		&=
		\begin{bmatrix}
			0\\0\\1
		\end{bmatrix}\\
		&=
		1
		\begin{bmatrix}
			0\\0\\1
		\end{bmatrix}\\
		&=1e_3.
	\end{aligned}
	\]
	
	Therefore,
	
	\[
	\boxed{
		\lambda_3=1,\qquad
		e_3=
		\begin{bmatrix}
			0\\0\\1
		\end{bmatrix}
	}.
	\]
	
	These eigenvectors are already normalized because
	
	\[
	e_1'e_1=e_2'e_2=e_3'e_3=1.
	\]
	
	Thus,
	
	\[
	\boxed{
		\begin{aligned}
			\lambda_1&=4,
			&
			e_1&=
			\begin{bmatrix}
				1\\0\\0
			\end{bmatrix},
			\\[4pt]
			\lambda_2&=9,
			&
			e_2&=
			\begin{bmatrix}
				0\\1\\0
			\end{bmatrix},
			\\[4pt]
			\lambda_3&=1,
			&
			e_3&=
			\begin{bmatrix}
				0\\0\\1
			\end{bmatrix}.
		\end{aligned}
	}
	\]
	
	\subsection*{(c) Eigenvalues and Eigenvectors of $\Sigma^{-1}$}
	
	From part (a),
	
	\[
	\Sigma^{-1}
	=
	\begin{bmatrix}
		1/4&0&0\\
		0&1/9&0\\
		0&0&1
	\end{bmatrix}.
	\]
	
	Again, $\Sigma^{-1}$ is diagonal and therefore triangular.
	Its eigenvalues are its diagonal elements:
	
	\[
	\boxed{
		\lambda_1=\frac14,\qquad
		\lambda_2=\frac19,\qquad
		\lambda_3=1
	}.
	\]
	
	The corresponding eigenvectors remain the same standard basis
	vectors:
	
	\[
	\boxed{
		e_1=
		\begin{bmatrix}
			1\\0\\0
		\end{bmatrix},
		\qquad
		e_2=
		\begin{bmatrix}
			0\\1\\0
		\end{bmatrix},
		\qquad
		e_3=
		\begin{bmatrix}
			0\\0\\1
		\end{bmatrix}
	}.
	\]
	
	For example,
	
	\[
	\begin{aligned}
		\Sigma^{-1}e_1
		&=
		\begin{bmatrix}
			1/4&0&0\\
			0&1/9&0\\
			0&0&1
		\end{bmatrix}
		\begin{bmatrix}
			1\\0\\0
		\end{bmatrix}\\
		&=
		\begin{bmatrix}
			1/4\\0\\0
		\end{bmatrix}\\
		&=
		\frac14e_1.
	\end{aligned}
	\]
	
	Thus, the eigenvalue changes from $4$ to $1/4$ while the
	eigenvector remains $e_1$.
	
	In general, if
	
	\[
	Ax=\lambda x,
	\]
	
	then, provided $\lambda\neq0$,
	
	\[
	A^{-1}x=\frac{1}{\lambda}x.
	\]
	
	Therefore,
	
	\[
	\boxed{
		\lambda_i(A^{-1})
		=
		\frac{1}{\lambda_i(A)}
	}
	\]
	
	while the corresponding eigenvectors remain unchanged.
	
	For this problem,
	
	\[
	4\longrightarrow\frac14,
	\qquad
	9\longrightarrow\frac19,
	\qquad
	1\longrightarrow1.
	\]
	
	\subsection*{Properties Used}
	
	The main properties used in this problem are
	
	\[
	\boxed{
		\text{Diagonal matrix}
		\Rightarrow
		\text{triangular matrix}
	}
	\]
	
	\[
	\boxed{
		\text{Triangular matrix}
		\Rightarrow
		\text{eigenvalues are diagonal elements}
	}
	\]
	
	\[
	\boxed{
		\text{Diagonal matrix}
		\Rightarrow
		\text{standard basis vectors give the corresponding eigenvectors}
	}
	\]
	
	because
	
	\[
	\Sigma
	\begin{bmatrix}
		1\\0\\0
	\end{bmatrix}
	=
	4
	\begin{bmatrix}
		1\\0\\0
	\end{bmatrix},
	\qquad
	\Sigma
	\begin{bmatrix}
		0\\1\\0
	\end{bmatrix}
	=
	9
	\begin{bmatrix}
		0\\1\\0
	\end{bmatrix},
	\]
	
	and
	
	\[
	\Sigma
	\begin{bmatrix}
		0\\0\\1
	\end{bmatrix}
	=
	1
	\begin{bmatrix}
		0\\0\\1
	\end{bmatrix}.
	\]
	
	Also,
	
	\[
	\boxed{
		\Sigma\Sigma^{-1}=I
		\Rightarrow
		\text{reciprocal diagonal elements for a diagonal inverse}
	}
	\]
	
	and
	
	\[
	\boxed{
		Ax=\lambda x
		\Rightarrow
		A^{-1}x=\frac{1}{\lambda}x
	}.
	\]
	
	Thus, 2.24 can be solved primarily by recognizing these
	matrix properties rather than repeatedly calculating a
	characteristic determinant.

	\section*{Problem 2.30}
	
	Given
	\[
	X=
	\begin{bmatrix}
		X_1\\
		X_2\\
		X_3\\
		X_4
	\end{bmatrix}
	=
	\begin{bmatrix}
		X^{(1)}\\
		X^{(2)}
	\end{bmatrix},
	\qquad
	X^{(1)}=
	\begin{bmatrix}
		X_1\\
		X_2
	\end{bmatrix},
	\qquad
	X^{(2)}=
	\begin{bmatrix}
		X_3\\
		X_4
	\end{bmatrix}.
	\]
	
	The mean vector and covariance matrix are
	\[
	\mu_X=
	\begin{bmatrix}
		4\\
		3\\
		2\\
		1
	\end{bmatrix},
	\qquad
	\Sigma_X=
	\begin{bmatrix}
		3&0&2&2\\
		0&1&1&0\\
		2&1&9&-2\\
		2&0&-2&4
	\end{bmatrix}.
	\]
	
	The covariance matrix is defined by
	\[
	\boxed{
		\Sigma_X
		=
		\operatorname{Cov}(X)
		=
		E\left[(X-\mu_X)(X-\mu_X)'\right]
	}
	\]
	
	where $E[\cdot]$ is the expected-value operator and $\mu_X=E(X)$.
	
	% -----------------------------------------------------
	\subsection*{(c) Find $\operatorname{Cov}(X^{(1)})$}
	
	Since
	\[
	X^{(1)}=
	\begin{bmatrix}
		X_1\\
		X_2
	\end{bmatrix},
	\]
	we take the covariance information for $X_1$ and $X_2$
	from the first two rows and columns of $\Sigma_X$.
	
	Thus,
	\[
	\operatorname{Cov}(X^{(1)})
	=
	\begin{bmatrix}
		3&0\\
		0&1
	\end{bmatrix}.
	\]
	
	The entries represent
	\[
	3=\operatorname{Var}(X_1),\qquad
	1=\operatorname{Var}(X_2),
	\]
	and
	\[
	0=\operatorname{Cov}(X_1,X_2).
	\]
	
	Therefore,
	\[
	\boxed{
		\operatorname{Cov}(X^{(1)})
		=
		\begin{bmatrix}
			3&0\\
			0&1
		\end{bmatrix}
	}.
	\]
	
	% -----------------------------------------------------
	\subsection*{(d) Find $\operatorname{Cov}(AX^{(1)})$}
	
	Given
	\[
	A=
	\begin{bmatrix}
		1&2
	\end{bmatrix},
	\]
	we first recognize that
	\[
	AX^{(1)}
	=
	\begin{bmatrix}
		1&2
	\end{bmatrix}
	\begin{bmatrix}
		X_1\\
		X_2
	\end{bmatrix}
	=
	X_1+2X_2.
	\]
	
	For a linear transformation $Y=AX$, covariance transforms as
	\[
	\boxed{
		\operatorname{Cov}(Y)
		=
		A\operatorname{Cov}(X)A'
	}
	\]
	
	because this follows directly from
	\[
	\operatorname{Cov}(Y)
	=
	E[(Y-\mu_Y)(Y-\mu_Y)']
	\]
	and
	\[
	Y-\mu_Y=A(X-\mu_X).
	\]
	
	Therefore,
	\[
	\operatorname{Cov}(AX^{(1)})
	=
	A\operatorname{Cov}(X^{(1)})A'.
	\]
	
	Since
	\[
	A'=
	\begin{bmatrix}
		1\\
		2
	\end{bmatrix},
	\]
	we get
	\[
	\begin{aligned}
		\operatorname{Cov}(AX^{(1)})
		&=
		\begin{bmatrix}
			1&2
		\end{bmatrix}
		\begin{bmatrix}
			3&0\\
			0&1
		\end{bmatrix}
		\begin{bmatrix}
			1\\
			2
		\end{bmatrix}\\
		&=
		\begin{bmatrix}
			3&2
		\end{bmatrix}
		\begin{bmatrix}
			1\\
			2
		\end{bmatrix}\\
		&=3+4\\
		&=7.
	\end{aligned}
	\]
	
	Thus,
	\[
	\boxed{
		\operatorname{Cov}(AX^{(1)})=7
	}
	\]
	
	Since $AX^{(1)}$ is only one variable,
	this covariance is actually its variance:
	\[
	\boxed{
		\operatorname{Cov}(AX^{(1)})
		=
		\operatorname{Var}(X_1+2X_2)=7.
	}
	\]
	
	% -----------------------------------------------------
	\subsection*{(e) Find $E(X^{(2)})$}
	
	Since
	\[
	X^{(2)}
	=
	\begin{bmatrix}
		X_3\\
		X_4
	\end{bmatrix},
	\]
	we take the last two entries of $\mu_X$:
	\[
	\boxed{
		E(X^{(2)})
		=
		\begin{bmatrix}
			2\\
			1
		\end{bmatrix}
	}.
	\]
	
	% -----------------------------------------------------
	\subsection*{(f) Find $E(BX^{(2)})$}
	
	Given
	\[
	B=
	\begin{bmatrix}
		1&-2\\
		2&-1
	\end{bmatrix},
	\]
	linearity of expectation gives
	\[
	\boxed{
		E(BX^{(2)})=BE(X^{(2)}).
	}
	\]
	
	Therefore,
	\[
	\begin{aligned}
		E(BX^{(2)})
		&=
		\begin{bmatrix}
			1&-2\\
			2&-1
		\end{bmatrix}
		\begin{bmatrix}
			2\\
			1
		\end{bmatrix}\\
		&=
		\begin{bmatrix}
			(1)(2)+(-2)(1)\\
			(2)(2)+(-1)(1)
		\end{bmatrix}\\
		&=
		\begin{bmatrix}
			0\\
			3
		\end{bmatrix}.
	\end{aligned}
	\]
	
	Thus,
	\[
	\boxed{
		E(BX^{(2)})
		=
		\begin{bmatrix}
			0\\
			3
		\end{bmatrix}
	}.
	\]
	
	% -----------------------------------------------------
	\subsection*{(g) Find $\operatorname{Cov}(X^{(2)})$}
	
	Since
	\[
	X^{(2)}
	=
	\begin{bmatrix}
		X_3\\
		X_4
	\end{bmatrix},
	\]
	we take the bottom-right block of $\Sigma_X$:
	
	\[
	\Sigma_{22}
	=
	\begin{bmatrix}
		9&-2\\
		-2&4
	\end{bmatrix}.
	\]
	
	Therefore,
	\[
	\boxed{
		\operatorname{Cov}(X^{(2)})
		=
		\begin{bmatrix}
			9&-2\\
			-2&4
		\end{bmatrix}
	}.
	\]
	
	Here,
	\[
	9=\operatorname{Var}(X_3),\qquad
	4=\operatorname{Var}(X_4),
	\]
	while
	\[
	-2=\operatorname{Cov}(X_3,X_4).
	\]
	
	The nonzero off-diagonal value means $X_3$ and $X_4$
	have covariance $-2$.
	
	% -----------------------------------------------------
	\subsection*{(h) Find $\operatorname{Cov}(BX^{(2)})$}
	
	Using the covariance transformation rule,
	\[
	\operatorname{Cov}(BX^{(2)})
	=
	B\operatorname{Cov}(X^{(2)})B'.
	\]
	
	Since
	\[
	B'=
	\begin{bmatrix}
		1&2\\
		-2&-1
	\end{bmatrix},
	\]
	we obtain
	\[
	\begin{aligned}
		\operatorname{Cov}(BX^{(2)})
		&=
		\begin{bmatrix}
			1&-2\\
			2&-1
		\end{bmatrix}
		\begin{bmatrix}
			9&-2\\
			-2&4
		\end{bmatrix}
		\begin{bmatrix}
			1&2\\
			-2&-1
		\end{bmatrix}\\
		&=
		\begin{bmatrix}
			13&-10\\
			20&-8
		\end{bmatrix}
		\begin{bmatrix}
			1&2\\
			-2&-1
		\end{bmatrix}\\
		&=
		\begin{bmatrix}
			33&36\\
			36&48
		\end{bmatrix}.
	\end{aligned}
	\]
	
	Therefore,
	\[
	\boxed{
		\operatorname{Cov}(BX^{(2)})
		=
		\begin{bmatrix}
			33&36\\
			36&48
		\end{bmatrix}
	}.
	\]
	
	The transformed vector is
	\[
	BX^{(2)}
	=
	\begin{bmatrix}
		X_3-2X_4\\
		2X_3-X_4
	\end{bmatrix}.
	\]
	
	Thus,
	\[
	33=\operatorname{Var}(X_3-2X_4),
	\]
	\[
	48=\operatorname{Var}(2X_3-X_4),
	\]
	and
	\[
	36=
	\operatorname{Cov}(X_3-2X_4,\;2X_3-X_4).
	\]
	
	% -----------------------------------------------------
	\subsection*{(i) Find $\operatorname{Cov}(X^{(1)},X^{(2)})$}
	
	Partition the covariance matrix according to
	\[
	X=
	\begin{bmatrix}
		X^{(1)}\\
		X^{(2)}
	\end{bmatrix}.
	\]
	
	Thus,
	\[
	\Sigma_X
	=
	\begin{bmatrix}
		\Sigma_{11}&\Sigma_{12}\\
		\Sigma_{21}&\Sigma_{22}
	\end{bmatrix}.
	\]
	
	The covariance between the first and second partitions
	is the top-right block:
	
	\[
	\boxed{
		\operatorname{Cov}(X^{(1)},X^{(2)})
		=
		\Sigma_{12}
		=
		\begin{bmatrix}
			2&2\\
			1&0
		\end{bmatrix}
	}.
	\]
	
	The rows correspond to $X_1,X_2$ and the columns
	correspond to $X_3,X_4$:
	
	\[
	\begin{array}{c|cc}
		& X_3 & X_4\\
		\hline
		X_1 & 2 & 2\\
		X_2 & 1 & 0
	\end{array}
	\]
	
	Thus,
	\[
	2=\operatorname{Cov}(X_1,X_3),
	\qquad
	2=\operatorname{Cov}(X_1,X_4),
	\]
	\[
	1=\operatorname{Cov}(X_2,X_3),
	\qquad
	0=\operatorname{Cov}(X_2,X_4).
	\]
	
	Because covariance is symmetric,
	\[
	\boxed{
		\operatorname{Cov}(X^{(2)},X^{(1)})
		=
		\operatorname{Cov}(X^{(1)},X^{(2)})'
	}
	\]
	
	and therefore
	
	\[
	\operatorname{Cov}(X^{(2)},X^{(1)})
	=
	\begin{bmatrix}
		2&1\\
		2&0
	\end{bmatrix}.
	\]
	
	% -----------------------------------------------------
	\subsection*{(j) Find $\operatorname{Cov}(AX^{(1)},BX^{(2)})$}
	
	For two linear transformations,
	\[
	\boxed{
		\operatorname{Cov}(AX^{(1)},BX^{(2)})
		=
		A\operatorname{Cov}(X^{(1)},X^{(2)})B'
	}.
	\]
	
	Using
	\[
	A=
	\begin{bmatrix}
		1&2
	\end{bmatrix},
	\qquad
	\operatorname{Cov}(X^{(1)},X^{(2)})
	=
	\begin{bmatrix}
		2&2\\
		1&0
	\end{bmatrix},
	\]
	and
	\[
	B'=
	\begin{bmatrix}
		1&2\\
		-2&-1
	\end{bmatrix},
	\]
	
	we get
	
	\[
	\begin{aligned}
		\operatorname{Cov}(AX^{(1)},BX^{(2)})
		&=
		\begin{bmatrix}
			1&2
		\end{bmatrix}
		\begin{bmatrix}
			2&2\\
			1&0
		\end{bmatrix}
		\begin{bmatrix}
			1&2\\
			-2&-1
		\end{bmatrix}\\
		&=
		\begin{bmatrix}
			4&2
		\end{bmatrix}
		\begin{bmatrix}
			1&2\\
			-2&-1
		\end{bmatrix}\\
		&=
		\begin{bmatrix}
			0&6
		\end{bmatrix}.
	\end{aligned}
	\]
	
	Therefore,
	
	\[
	\boxed{
		\operatorname{Cov}(AX^{(1)},BX^{(2)})
		=
		\begin{bmatrix}
			0&6
		\end{bmatrix}
	}.
	\]
	
	Here,
	\[
	AX^{(1)}=X_1+2X_2
	\]
	
	is one new variable, while
	
	\[
	BX^{(2)}
	=
	\begin{bmatrix}
		X_3-2X_4\\
		2X_3-X_4
	\end{bmatrix}
	\]
	
	contains two new variables.
	
	Thus,
	
	\[
	0=
	\operatorname{Cov}
	(X_1+2X_2,\;X_3-2X_4),
	\]
	
	and
	
	\[
	6=
	\operatorname{Cov}
	(X_1+2X_2,\;2X_3-X_4).
	\]
	
	\subsection*{Problem 2.41}
	
	Given
	
	\[
	\mu_X=E(X)=
	\begin{bmatrix}
		3\\
		2\\
		-2\\
		0
	\end{bmatrix},
	\qquad
	\Sigma_X=\operatorname{Cov}(X)
	=
	\begin{bmatrix}
		3&0&0&0\\
		0&3&0&0\\
		0&0&3&0\\
		0&0&0&3
	\end{bmatrix}
	=3I.
	\]
	
	Let
	
	\[
	A=
	\begin{bmatrix}
		1&-1&0&0\\
		1&1&-2&0\\
		1&1&1&-3
	\end{bmatrix}.
	\]
	
	Since $Y=AX$, the rows of $A$ create the three linear
	combinations
	
	\[
	Y_1=X_1-X_2,
	\]
	
	\[
	Y_2=X_1+X_2-2X_3,
	\]
	
	\[
	Y_3=X_1+X_2+X_3-3X_4.
	\]
	
	% -----------------------------------------------------
	
	\subsubsection*{(a) Find $E(AX)$}
	
	For a linear transformation,
	
	\[
	\boxed{
		E(AX)=AE(X)=A\mu_X
	}.
	\]
	
	Therefore,
	
	\[
	\begin{aligned}
		E(AX)
		&=
		\begin{bmatrix}
			1&-1&0&0\\
			1&1&-2&0\\
			1&1&1&-3
		\end{bmatrix}
		\begin{bmatrix}
			3\\
			2\\
			-2\\
			0
		\end{bmatrix}\\
		&=
		\begin{bmatrix}
			(1)(3)+(-1)(2)+(0)(-2)+(0)(0)\\
			(1)(3)+(1)(2)+(-2)(-2)+(0)(0)\\
			(1)(3)+(1)(2)+(1)(-2)+(-3)(0)
		\end{bmatrix}\\
		&=
		\begin{bmatrix}
			3-2\\
			3+2+4\\
			3+2-2
		\end{bmatrix}\\
		&=
		\boxed{
			\begin{bmatrix}
				1\\
				9\\
				3
			\end{bmatrix}
		}.
	\end{aligned}
	\]
	
	% -----------------------------------------------------
	
	\subsubsection*{(b) Find $\operatorname{Cov}(AX)$}
	
	For a linear transformation,
	
	\[
	\boxed{
		\operatorname{Cov}(AX)=A\Sigma_XA'
	}.
	\]
	
	Since
	
	\[
	\Sigma_X=3I,
	\]
	
	we can simplify:
	
	\[
	\begin{aligned}
		\operatorname{Cov}(AX)
		&=A(3I)A'\\
		&=3AIA'\\
		&=3AA'.
	\end{aligned}
	\]
	
	The transpose of $A$ is
	
	\[
	A'=
	\begin{bmatrix}
		1&1&1\\
		-1&1&1\\
		0&-2&1\\
		0&0&-3
	\end{bmatrix}.
	\]
	
	Let the rows of $A$ be
	
	\[
	r_1=(1,-1,0,0),
	\]
	
	\[
	r_2=(1,1,-2,0),
	\]
	
	\[
	r_3=(1,1,1,-3).
	\]
	
	Then $AA'$ contains the dot products of the rows:
	
	\[
	AA'
	=
	\begin{bmatrix}
		r_1r_1'&r_1r_2'&r_1r_3'\\
		r_2r_1'&r_2r_2'&r_2r_3'\\
		r_3r_1'&r_3r_2'&r_3r_3'
	\end{bmatrix}.
	\]
	
	Calculate the unique entries:
	
	\[
	\begin{aligned}
		r_1r_1'
		&=(1)^2+(-1)^2
		=2,\\
		r_1r_2'
		&=(1)(1)+(-1)(1)
		=0,\\
		r_1r_3'
		&=(1)(1)+(-1)(1)
		=0,\\
		r_2r_2'
		&=(1)^2+(1)^2+(-2)^2
		=6,\\
		r_2r_3'
		&=(1)(1)+(1)(1)+(-2)(1)
		=0,\\
		r_3r_3'
		&=(1)^2+(1)^2+(1)^2+(-3)^2
		=12.
	\end{aligned}
	\]
	
	Because $AA'$ is symmetric,
	
	\[
	AA'
	=
	\begin{bmatrix}
		2&0&0\\
		0&6&0\\
		0&0&12
	\end{bmatrix}.
	\]
	
	Thus,
	
	\[
	\begin{aligned}
		\operatorname{Cov}(AX)
		&=3AA'\\
		&=
		3
		\begin{bmatrix}
			2&0&0\\
			0&6&0\\
			0&0&12
		\end{bmatrix}\\
		&=
		\boxed{
			\begin{bmatrix}
				6&0&0\\
				0&18&0\\
				0&0&36
			\end{bmatrix}
		}.
	\end{aligned}
	\]
	
	The diagonal entries are the variances:
	
	\[
	\operatorname{Var}(Y_1)=6,
	\qquad
	\operatorname{Var}(Y_2)=18,
	\qquad
	\operatorname{Var}(Y_3)=36.
	\]
	
	The off-diagonal entries are the covariances:
	
	\[
	\operatorname{Cov}(Y_1,Y_2)
	=
	\operatorname{Cov}(Y_1,Y_3)
	=
	\operatorname{Cov}(Y_2,Y_3)
	=0.
	\]
	
	% -----------------------------------------------------
	
	\subsubsection*{(c) Which pairs of linear combinations have zero covariances?}
	
	Since
	
	\[
	\operatorname{Cov}(AX)
	=
	\begin{bmatrix}
		6&0&0\\
		0&18&0\\
		0&0&36
	\end{bmatrix},
	\]
	
	all off-diagonal entries are zero.
	
	Therefore, all three pairs have zero covariance:
	
	\[
	\boxed{
		(Y_1,Y_2),\qquad
		(Y_1,Y_3),\qquad
		(Y_2,Y_3)
	}.
	\]
	
	Using the original linear combinations:
	
	\[
	\boxed{
		(X_1-X_2,\;X_1+X_2-2X_3)
	}
	\]
	
	\[
	\boxed{
		(X_1-X_2,\;X_1+X_2+X_3-3X_4)
	}
	\]
	
	\[
	\boxed{
		(X_1+X_2-2X_3,\;
		X_1+X_2+X_3-3X_4)
	}.
	\]
	
	All three pairs have covariance $0$.
	
	% -----------------------------------------------------
	
	\subsubsection*{Key Ideas}
	
	For a linear transformation $Y=AX$,
	
	\[
	\boxed{
		E(Y)=AE(X)
	}
	\]
	
	and
	
	\[
	\boxed{
		\operatorname{Cov}(Y)=A\Sigma_XA'
	}.
	\]
	
	Here,
	
	\[
	\Sigma_X=3I,
	\]
	
	so
	
	\[
	\boxed{
		\operatorname{Cov}(AX)=3AA'
	}.
	\]
	
	The rows of $A$ determine the new linear combinations.
	Since $A$ has three rows, $AX$ contains three new variables,
	so its covariance matrix is $3\times3$.
	
	The diagonal entries give the variances, while the
	off-diagonal entries give the covariances.
	
\end{document}